% Общая реализация игрового поля «шесть пешек» (Hexapawn).
%
% Файл НЕ самостоятельный: подключается через % Общая реализация игрового поля «шесть пешек» (Hexapawn).
%
% Файл НЕ самостоятельный: подключается через % Общая реализация игрового поля «шесть пешек» (Hexapawn).
%
% Файл НЕ самостоятельный: подключается через % Общая реализация игрового поля «шесть пешек» (Hexapawn).
%
% Файл НЕ самостоятельный: подключается через \input{board.tex} из rules.tex,
% examples.tex, cards.tex, field.tex и guide.tex. rating.tex не подключает:
% ему нужны из палитры только цвета, поэтому он берёт colors.tex напрямую.
% Содержит только определения — \documentclass здесь быть не должно.
%
% Квадраты нумеруются слева направо и снизу вверх, как на настоящей доске:
%   7 8 9   <- ряд 3
%   4 5 6   <- ряд 2
%   1 2 3   <- ряд 1
% То есть a1=1, b1=2, c1=3, a2=4, …, c3=9.
% Ряд 1 внизу: белые стоят в ряду 1, чёрные в ряду 3.
% Подписи совпадают с подписями на игровом поле.
%
% Что где живёт:
%   \sqcx, \sqcy      — геометрия квадрата по его номеру;
%   \boardgrid        — тонкая сетка (карточки, игровое поле);
%   \boardgridbold    — сетка одной толщины с границей (пустое поле);
%   \boardlabels      — подписи a/b/c и 3/2/1 по краям доски 5,4 см;
%   \boardtitle       — заголовок над доской 5,4 см;
%   \drawpawns        — пешки, заданные номерами квадратов;
%   \board, \boardmoves — небольшие иллюстрации к тексту (доска 1,9 см);
%   \card, \hexcard, \cardfoot — карточка целиком (доска 5,4 см).
%
% Масштаб доски задаётся стилем TikZ, который выбирает вызывающий документ:
% rulesheet — 1,9 см (иллюстрации в тексте),
% cardsheet — 5,4 см (карточки).
% =====================================================================================================================

\RequirePackage{tikz}
\usetikzlibrary{arrows.meta}

% Палитра вынесена в colors.tex, см. там же список цветов и что ими красится.
\input{colors.tex}


% =====================================================================================================================
% Геометрия квадрата по его номеру
% =====================================================================================================================
\newcommand{\sqcx}[1]{\pgfmathsetmacro{\qx}{mod(#1-1,3)+0.5}}       % центр по X
\newcommand{\sqcy}[1]{\pgfmathsetmacro{\qy}{floor((#1-1)/3)+0.5}}   % центр по Y

% =====================================================================================================================
% Сетка доски
% =====================================================================================================================
% Тонкая: внутренние линии светло-серые, граница толще. Карточки и поле.
\newcommand{\boardgrid}{%
  \foreach \gl in {1,2} {\draw[gridline,line width=0.4pt] (\gl,0) -- (\gl,3) (0,\gl) -- (3,\gl);}%
  \draw[ink,line width=1.2pt] (0,0) rectangle (3,3);%
}
% Для пустой доски: сетка той же толщины и цвета, что и граница.
\newcommand{\boardgridbold}{%
  \foreach \gl in {1,2} {\draw[ink,line width=1.2pt] (\gl,0) -- (\gl,3) (0,\gl) -- (3,\gl);}%
  \draw[ink,line width=1.2pt] (0,0) rectangle (3,3);%
}

% =====================================================================================================================
% Подписи и заголовок для доски 5,4 см. #1 --- стиль подписи, #2 --- заголовка.
% Смещение 0,70 см пересчитано в единицы доски делением на размер клетки.
% =====================================================================================================================
\newcommand{\boardlabels}[1]{%
  \foreach \cl/\lab in {0/a,1/b,2/c} {%
    \node[#1] at (\cl+0.5,{-0.70/5.4}) {\lab};
    \node[#1] at (\cl+0.5,{3.0+0.70/5.4}) {\lab};}%
  \foreach \rw/\lab in {0/3,1/2,2/1} {%
    \node[#1] at ({-0.70/5.4},{2.5-\rw}) {\lab};
    \node[#1] at ({3.0+0.70/5.4},{2.5-\rw}) {\lab};}%
}
\newcommand{\boardtitle}[2]{%
  \node[#2] at (1.5,{3.0+3.15/5.4}) {#1};%
}

% =====================================================================================================================
% Пешки по номерам квадратов. Радиус задан в единицах доски.
% =====================================================================================================================
\newcommand{\drawpawns}[2]{%
  \foreach \sqr in {#2} {\sqcx{\sqr}\sqcy{\sqr}\fill[pawn] (\qx,\qy) circle[radius=0.34];}%
  \foreach \sqr in {#1} {\sqcx{\sqr}\sqcy{\sqr}%
    \filldraw[fill=white,draw=pawn,line width=1.1pt] (\qx,\qy) circle[radius=0.34];}%
}

% =====================================================================================================================
% Небольшие иллюстрации к тексту правил и примеров (доска 1,9 см).
% Здесь пешки рисуются кружками-узлами, а не кругами: на маленьком размере так они лучше читаются.
% =====================================================================================================================
\tikzset{
  rulesheet/.style={
    x=1.9cm, y=1.9cm,
    line width=0.5pt, baseline=(current bounding box.center)},
  plainpawn/.style={circle, draw=pawn, line width=0.8pt, minimum size=7mm, inner sep=0pt},
  blpawn/.style={plainpawn, fill=pawn, draw=pawn},
  whpawn/.style={plainpawn, fill=white, draw=pawn},
  moveline/.style={-{Stealth[length=2.4mm,width=1.7mm]}, line width=0.9pt, draw=moveink},
}
\newcommand{\rulesheetgrid}{%
  \foreach \gl in {1,2} {\draw[ink!35] (\gl,0) -- (\gl,3) (0,\gl) -- (3,\gl);}%
  \draw[ink, line width=1pt] (0,0) rectangle (3,3);%
}
% \board{<белые>}{<чёрные>}{подписи: 1 да, 0 нет}
\newcommand{\board}[3]{%
  \begin{tikzpicture}[rulesheet]
    \rulesheetgrid
    \foreach \s in {#1} {\sqcx{\s}\sqcy{\s}\node[whpawn] at (\qx,\qy) {};}
    \foreach \s in {#2} {\sqcx{\s}\sqcy{\s}\node[blpawn] at (\qx,\qy) {};}
    \ifnum#3=1\relax
      \foreach \c/\l in {0/a,1/b,2/c} {%
        \node[font=\scriptsize\sffamily, text=ink!70] at (\c+0.5,{-0.34}) {\l};
        \node[font=\scriptsize\sffamily, text=ink!70] at (\c+0.5,{3.34}) {\l};}
      \foreach \r in {1,2,3} {%
        \node[font=\scriptsize\sffamily, text=ink!70] at ({-0.32},{\r-0.5}) {\r};
        \node[font=\scriptsize\sffamily, text=ink!70] at ({3.32},{\r-0.5}) {\r};}
    \fi
  \end{tikzpicture}%
}
% \boardmoves{<белые>}{<чёрные>}{<ходы: from/to, ...>}
\newcommand{\boardmoves}[3]{%
  \begin{tikzpicture}[rulesheet]
    \rulesheetgrid
    \foreach \s in {#1} {\sqcx{\s}\sqcy{\s}\node[whpawn] at (\qx,\qy) {};}
    \foreach \s in {#2} {\sqcx{\s}\sqcy{\s}\node[blpawn] at (\qx,\qy) {};}
    \foreach \f/\t in {#3} {%
      \sqcx{\f}\sqcy{\f}\pgfmathsetmacro{\fx}{\qx}\pgfmathsetmacro{\fy}{\qy}%
      \sqcx{\t}\sqcy{\t}%
      \draw[moveline] (\fx,\fy) -- (\qx,\qy);}
  \end{tikzpicture}%
}
% Крестик взятия: \take в тексте.
\newcommand{\take}{\ensuremath{\times}}

% =====================================================================================================================
% Карточка целиком (доска 5,4 см). Здесь пешки — круги, стрелки обрезаются
% по краям, чтобы пешки оставались видимыми.
% =====================================================================================================================
\tikzset{
  cardsheet/.style={x=5.4cm, y=5.4cm, line width=0.7pt, baseline=(current bounding box.center)},
  movetip/.style={-{Stealth[length=4mm,width=3.2mm]}, line width=5pt},
  squarelab/.style={font=\sffamily\bfseries\fontsize{36}{44}\selectfont, text=ink},
  movelab/.style={font=\sffamily\bfseries\fontsize{28}{34}\selectfont, text=ink},
  fieldlab/.style={font=\sffamily\bfseries\fontsize{36}{44}\selectfont,text=ink},
  fieldtitle/.style={font=\sffamily\fontsize{28}{34}\selectfont\bfseries,text=ink},
}
% Насколько конец стрелки не доходит до центра квадрата назначения:
% 0.13 для тихого хода, 0.54 для взятия.
\newcommand{\quietstop}{0.13}
\newcommand{\capturestop}{0.54}
% \card{<белые>}{<чёрные>}{<ходы>}{<подпись>}
% ходы перечисляются как <откуда>/<куда>/<цвет>, например 3/6/arrow1 по одному на каждый допустимый ход из этой позиции.
% Стрелка идёт от пешки к пешке, но обрезается по краям:
%   начало выходит из-под своей пешки,
%   конец уходит под пешку назначения (тихий ход) или останавливается перед чужой (взятие).
% Так пешки всегда остаются видимыми.
\newcommand{\card}[4]{%
  \begin{tikzpicture}[cardsheet]
    \boardgrid
    \foreach \mvfrom/\mvto/\mvcolor in {#3} {%
      \global\let\stopat\quietstop
      \foreach \sqr in {#1} {\ifnum\sqr=\mvto\relax\global\let\stopat\capturestop\fi}%
      \foreach \sqr in {#2} {\ifnum\sqr=\mvto\relax\global\let\stopat\capturestop\fi}%
      \sqcx{\mvfrom}\sqcy{\mvfrom}\pgfmathsetmacro{\fx}{\qx}\pgfmathsetmacro{\fy}{\qy}%
      \sqcx{\mvto}\sqcy{\mvto}%
      \pgfmathsetmacro{\ux}{\qx-\fx}\pgfmathsetmacro{\uy}{\qy-\fy}%
      \pgfmathsetmacro{\ulen}{sqrt(\ux*\ux+\uy*\uy)}%
      \pgfmathsetmacro{\sx}{round((\fx+0.40*\ux/\ulen)*100)/100}%
      \pgfmathsetmacro{\sy}{round((\fy+0.40*\uy/\ulen)*100)/100}%
      \pgfmathsetmacro{\ex}{round((\qx-\stopat*\ux/\ulen)*100)/100}%
      \pgfmathsetmacro{\ey}{round((\qy-\stopat*\uy/\ulen)*100)/100}%
      \draw[movetip,\mvcolor] (\sx,\sy) -- (\ex,\ey);%
    }%
    \drawpawns{#1}{#2}%
    \boardlabels{squarelab}%
    \boardtitle{#4}{movelab}%
  \end{tikzpicture}%
}
% Карточка целиком: доска по центру листа.
\newcommand{\hexcard}[4]{%
  \vspace*{\fill}%
  \begin{center}\card{#1}{#2}{#3}{#4}\end{center}%
  \vspace*{\fill}%
}
% Колонтитул: «страница/всего» слева, «карточка/колода» справа.
% Для страницы начальной расстановки номер в колоде не печатается.
\newcommand{\cardfoot}[2]{%
  \noindent\hbox to \textwidth{\color{black!30}\sffamily\fontsize{9}{11}\selectfont
    #1\if\relax\detokenize{#2}\relax\else\hfil #2\fi}%
  \vspace*{1.4cm}%
}
 из rules.tex,
% examples.tex, cards.tex, field.tex и guide.tex. rating.tex не подключает:
% ему нужны из палитры только цвета, поэтому он берёт colors.tex напрямую.
% Содержит только определения — \documentclass здесь быть не должно.
%
% Квадраты нумеруются слева направо и снизу вверх, как на настоящей доске:
%   7 8 9   <- ряд 3
%   4 5 6   <- ряд 2
%   1 2 3   <- ряд 1
% То есть a1=1, b1=2, c1=3, a2=4, …, c3=9.
% Ряд 1 внизу: белые стоят в ряду 1, чёрные в ряду 3.
% Подписи совпадают с подписями на игровом поле.
%
% Что где живёт:
%   \sqcx, \sqcy      — геометрия квадрата по его номеру;
%   \boardgrid        — тонкая сетка (карточки, игровое поле);
%   \boardgridbold    — сетка одной толщины с границей (пустое поле);
%   \boardlabels      — подписи a/b/c и 3/2/1 по краям доски 5,4 см;
%   \boardtitle       — заголовок над доской 5,4 см;
%   \drawpawns        — пешки, заданные номерами квадратов;
%   \board, \boardmoves — небольшие иллюстрации к тексту (доска 1,9 см);
%   \card, \hexcard, \cardfoot — карточка целиком (доска 5,4 см).
%
% Масштаб доски задаётся стилем TikZ, который выбирает вызывающий документ:
% rulesheet — 1,9 см (иллюстрации в тексте),
% cardsheet — 5,4 см (карточки).
% =====================================================================================================================

\RequirePackage{tikz}
\usetikzlibrary{arrows.meta}

% Палитра вынесена в colors.tex, см. там же список цветов и что ими красится.
% Общая палитра цветов для всего комплекта.
% Меняйте значения HEX и комментарии, имена менять не рекомендуется.
% Стрелки arrow1…arrow4 — условные названия, не переименовывать.

\RequirePackage{xcolor}

% Основная палитра для печати игровой доски и расстановки пешек
\definecolor{ink}{HTML}{26262A}       % текст, рамка доски, подписи клеток, оси в rating
\definecolor{pawn}{HTML}{000000}      % цвет пешек чёрных и контура пешек белых
\definecolor{gridline}{HTML}{B4B4BC}  % внутренние линии сетки
% Отдельный светлый цвет для стрелок на иллюстрациях: такая стрелка должна
% быть светлее чёрной пешки, иначе при взятии (например a2xb3) кончик
% стрелки пропадает на пешке. Поэтому у moveline свой цвет, а не ink.
\definecolor{moveink}{HTML}{B4B4BC}   % стрелки ходов в иллюстрациях и примерах

% Цвета стрелок: каждой стрелке на карточке соответствует конфета того же цвета.
\definecolor{arrow1}{HTML}{00B050}    % зелёный
\definecolor{arrow2}{HTML}{7030A0}    % фиолетовый
\definecolor{arrow3}{HTML}{FF0000}    % красный
\definecolor{arrow4}{HTML}{FFFF00}    % жёлтый



% =====================================================================================================================
% Геометрия квадрата по его номеру
% =====================================================================================================================
\newcommand{\sqcx}[1]{\pgfmathsetmacro{\qx}{mod(#1-1,3)+0.5}}       % центр по X
\newcommand{\sqcy}[1]{\pgfmathsetmacro{\qy}{floor((#1-1)/3)+0.5}}   % центр по Y

% =====================================================================================================================
% Сетка доски
% =====================================================================================================================
% Тонкая: внутренние линии светло-серые, граница толще. Карточки и поле.
\newcommand{\boardgrid}{%
  \foreach \gl in {1,2} {\draw[gridline,line width=0.4pt] (\gl,0) -- (\gl,3) (0,\gl) -- (3,\gl);}%
  \draw[ink,line width=1.2pt] (0,0) rectangle (3,3);%
}
% Для пустой доски: сетка той же толщины и цвета, что и граница.
\newcommand{\boardgridbold}{%
  \foreach \gl in {1,2} {\draw[ink,line width=1.2pt] (\gl,0) -- (\gl,3) (0,\gl) -- (3,\gl);}%
  \draw[ink,line width=1.2pt] (0,0) rectangle (3,3);%
}

% =====================================================================================================================
% Подписи и заголовок для доски 5,4 см. #1 --- стиль подписи, #2 --- заголовка.
% Смещение 0,70 см пересчитано в единицы доски делением на размер клетки.
% =====================================================================================================================
\newcommand{\boardlabels}[1]{%
  \foreach \cl/\lab in {0/a,1/b,2/c} {%
    \node[#1] at (\cl+0.5,{-0.70/5.4}) {\lab};
    \node[#1] at (\cl+0.5,{3.0+0.70/5.4}) {\lab};}%
  \foreach \rw/\lab in {0/3,1/2,2/1} {%
    \node[#1] at ({-0.70/5.4},{2.5-\rw}) {\lab};
    \node[#1] at ({3.0+0.70/5.4},{2.5-\rw}) {\lab};}%
}
\newcommand{\boardtitle}[2]{%
  \node[#2] at (1.5,{3.0+3.15/5.4}) {#1};%
}

% =====================================================================================================================
% Пешки по номерам квадратов. Радиус задан в единицах доски.
% =====================================================================================================================
\newcommand{\drawpawns}[2]{%
  \foreach \sqr in {#2} {\sqcx{\sqr}\sqcy{\sqr}\fill[pawn] (\qx,\qy) circle[radius=0.34];}%
  \foreach \sqr in {#1} {\sqcx{\sqr}\sqcy{\sqr}%
    \filldraw[fill=white,draw=pawn,line width=1.1pt] (\qx,\qy) circle[radius=0.34];}%
}

% =====================================================================================================================
% Небольшие иллюстрации к тексту правил и примеров (доска 1,9 см).
% Здесь пешки рисуются кружками-узлами, а не кругами: на маленьком размере так они лучше читаются.
% =====================================================================================================================
\tikzset{
  rulesheet/.style={
    x=1.9cm, y=1.9cm,
    line width=0.5pt, baseline=(current bounding box.center)},
  plainpawn/.style={circle, draw=pawn, line width=0.8pt, minimum size=7mm, inner sep=0pt},
  blpawn/.style={plainpawn, fill=pawn, draw=pawn},
  whpawn/.style={plainpawn, fill=white, draw=pawn},
  moveline/.style={-{Stealth[length=2.4mm,width=1.7mm]}, line width=0.9pt, draw=moveink},
}
\newcommand{\rulesheetgrid}{%
  \foreach \gl in {1,2} {\draw[ink!35] (\gl,0) -- (\gl,3) (0,\gl) -- (3,\gl);}%
  \draw[ink, line width=1pt] (0,0) rectangle (3,3);%
}
% \board{<белые>}{<чёрные>}{подписи: 1 да, 0 нет}
\newcommand{\board}[3]{%
  \begin{tikzpicture}[rulesheet]
    \rulesheetgrid
    \foreach \s in {#1} {\sqcx{\s}\sqcy{\s}\node[whpawn] at (\qx,\qy) {};}
    \foreach \s in {#2} {\sqcx{\s}\sqcy{\s}\node[blpawn] at (\qx,\qy) {};}
    \ifnum#3=1\relax
      \foreach \c/\l in {0/a,1/b,2/c} {%
        \node[font=\scriptsize\sffamily, text=ink!70] at (\c+0.5,{-0.34}) {\l};
        \node[font=\scriptsize\sffamily, text=ink!70] at (\c+0.5,{3.34}) {\l};}
      \foreach \r in {1,2,3} {%
        \node[font=\scriptsize\sffamily, text=ink!70] at ({-0.32},{\r-0.5}) {\r};
        \node[font=\scriptsize\sffamily, text=ink!70] at ({3.32},{\r-0.5}) {\r};}
    \fi
  \end{tikzpicture}%
}
% \boardmoves{<белые>}{<чёрные>}{<ходы: from/to, ...>}
\newcommand{\boardmoves}[3]{%
  \begin{tikzpicture}[rulesheet]
    \rulesheetgrid
    \foreach \s in {#1} {\sqcx{\s}\sqcy{\s}\node[whpawn] at (\qx,\qy) {};}
    \foreach \s in {#2} {\sqcx{\s}\sqcy{\s}\node[blpawn] at (\qx,\qy) {};}
    \foreach \f/\t in {#3} {%
      \sqcx{\f}\sqcy{\f}\pgfmathsetmacro{\fx}{\qx}\pgfmathsetmacro{\fy}{\qy}%
      \sqcx{\t}\sqcy{\t}%
      \draw[moveline] (\fx,\fy) -- (\qx,\qy);}
  \end{tikzpicture}%
}
% Крестик взятия: \take в тексте.
\newcommand{\take}{\ensuremath{\times}}

% =====================================================================================================================
% Карточка целиком (доска 5,4 см). Здесь пешки — круги, стрелки обрезаются
% по краям, чтобы пешки оставались видимыми.
% =====================================================================================================================
\tikzset{
  cardsheet/.style={x=5.4cm, y=5.4cm, line width=0.7pt, baseline=(current bounding box.center)},
  movetip/.style={-{Stealth[length=4mm,width=3.2mm]}, line width=5pt},
  squarelab/.style={font=\sffamily\bfseries\fontsize{36}{44}\selectfont, text=ink},
  movelab/.style={font=\sffamily\bfseries\fontsize{28}{34}\selectfont, text=ink},
  fieldlab/.style={font=\sffamily\bfseries\fontsize{36}{44}\selectfont,text=ink},
  fieldtitle/.style={font=\sffamily\fontsize{28}{34}\selectfont\bfseries,text=ink},
}
% Насколько конец стрелки не доходит до центра квадрата назначения:
% 0.13 для тихого хода, 0.54 для взятия.
\newcommand{\quietstop}{0.13}
\newcommand{\capturestop}{0.54}
% \card{<белые>}{<чёрные>}{<ходы>}{<подпись>}
% ходы перечисляются как <откуда>/<куда>/<цвет>, например 3/6/arrow1 по одному на каждый допустимый ход из этой позиции.
% Стрелка идёт от пешки к пешке, но обрезается по краям:
%   начало выходит из-под своей пешки,
%   конец уходит под пешку назначения (тихий ход) или останавливается перед чужой (взятие).
% Так пешки всегда остаются видимыми.
\newcommand{\card}[4]{%
  \begin{tikzpicture}[cardsheet]
    \boardgrid
    \foreach \mvfrom/\mvto/\mvcolor in {#3} {%
      \global\let\stopat\quietstop
      \foreach \sqr in {#1} {\ifnum\sqr=\mvto\relax\global\let\stopat\capturestop\fi}%
      \foreach \sqr in {#2} {\ifnum\sqr=\mvto\relax\global\let\stopat\capturestop\fi}%
      \sqcx{\mvfrom}\sqcy{\mvfrom}\pgfmathsetmacro{\fx}{\qx}\pgfmathsetmacro{\fy}{\qy}%
      \sqcx{\mvto}\sqcy{\mvto}%
      \pgfmathsetmacro{\ux}{\qx-\fx}\pgfmathsetmacro{\uy}{\qy-\fy}%
      \pgfmathsetmacro{\ulen}{sqrt(\ux*\ux+\uy*\uy)}%
      \pgfmathsetmacro{\sx}{round((\fx+0.40*\ux/\ulen)*100)/100}%
      \pgfmathsetmacro{\sy}{round((\fy+0.40*\uy/\ulen)*100)/100}%
      \pgfmathsetmacro{\ex}{round((\qx-\stopat*\ux/\ulen)*100)/100}%
      \pgfmathsetmacro{\ey}{round((\qy-\stopat*\uy/\ulen)*100)/100}%
      \draw[movetip,\mvcolor] (\sx,\sy) -- (\ex,\ey);%
    }%
    \drawpawns{#1}{#2}%
    \boardlabels{squarelab}%
    \boardtitle{#4}{movelab}%
  \end{tikzpicture}%
}
% Карточка целиком: доска по центру листа.
\newcommand{\hexcard}[4]{%
  \vspace*{\fill}%
  \begin{center}\card{#1}{#2}{#3}{#4}\end{center}%
  \vspace*{\fill}%
}
% Колонтитул: «страница/всего» слева, «карточка/колода» справа.
% Для страницы начальной расстановки номер в колоде не печатается.
\newcommand{\cardfoot}[2]{%
  \noindent\hbox to \textwidth{\color{black!30}\sffamily\fontsize{9}{11}\selectfont
    #1\if\relax\detokenize{#2}\relax\else\hfil #2\fi}%
  \vspace*{1.4cm}%
}
 из rules.tex,
% examples.tex, cards.tex, field.tex и guide.tex. rating.tex не подключает:
% ему нужны из палитры только цвета, поэтому он берёт colors.tex напрямую.
% Содержит только определения — \documentclass здесь быть не должно.
%
% Квадраты нумеруются слева направо и снизу вверх, как на настоящей доске:
%   7 8 9   <- ряд 3
%   4 5 6   <- ряд 2
%   1 2 3   <- ряд 1
% То есть a1=1, b1=2, c1=3, a2=4, …, c3=9.
% Ряд 1 внизу: белые стоят в ряду 1, чёрные в ряду 3.
% Подписи совпадают с подписями на игровом поле.
%
% Что где живёт:
%   \sqcx, \sqcy      — геометрия квадрата по его номеру;
%   \boardgrid        — тонкая сетка (карточки, игровое поле);
%   \boardgridbold    — сетка одной толщины с границей (пустое поле);
%   \boardlabels      — подписи a/b/c и 3/2/1 по краям доски 5,4 см;
%   \boardtitle       — заголовок над доской 5,4 см;
%   \drawpawns        — пешки, заданные номерами квадратов;
%   \board, \boardmoves — небольшие иллюстрации к тексту (доска 1,9 см);
%   \card, \hexcard, \cardfoot — карточка целиком (доска 5,4 см).
%
% Масштаб доски задаётся стилем TikZ, который выбирает вызывающий документ:
% rulesheet — 1,9 см (иллюстрации в тексте),
% cardsheet — 5,4 см (карточки).
% =====================================================================================================================

\RequirePackage{tikz}
\usetikzlibrary{arrows.meta}

% Палитра вынесена в colors.tex, см. там же список цветов и что ими красится.
% Общая палитра цветов для всего комплекта.
% Меняйте значения HEX и комментарии, имена менять не рекомендуется.
% Стрелки arrow1…arrow4 — условные названия, не переименовывать.

\RequirePackage{xcolor}

% Основная палитра для печати игровой доски и расстановки пешек
\definecolor{ink}{HTML}{26262A}       % текст, рамка доски, подписи клеток, оси в rating
\definecolor{pawn}{HTML}{000000}      % цвет пешек чёрных и контура пешек белых
\definecolor{gridline}{HTML}{B4B4BC}  % внутренние линии сетки
% Отдельный светлый цвет для стрелок на иллюстрациях: такая стрелка должна
% быть светлее чёрной пешки, иначе при взятии (например a2xb3) кончик
% стрелки пропадает на пешке. Поэтому у moveline свой цвет, а не ink.
\definecolor{moveink}{HTML}{B4B4BC}   % стрелки ходов в иллюстрациях и примерах

% Цвета стрелок: каждой стрелке на карточке соответствует конфета того же цвета.
\definecolor{arrow1}{HTML}{00B050}    % зелёный
\definecolor{arrow2}{HTML}{7030A0}    % фиолетовый
\definecolor{arrow3}{HTML}{FF0000}    % красный
\definecolor{arrow4}{HTML}{FFFF00}    % жёлтый



% =====================================================================================================================
% Геометрия квадрата по его номеру
% =====================================================================================================================
\newcommand{\sqcx}[1]{\pgfmathsetmacro{\qx}{mod(#1-1,3)+0.5}}       % центр по X
\newcommand{\sqcy}[1]{\pgfmathsetmacro{\qy}{floor((#1-1)/3)+0.5}}   % центр по Y

% =====================================================================================================================
% Сетка доски
% =====================================================================================================================
% Тонкая: внутренние линии светло-серые, граница толще. Карточки и поле.
\newcommand{\boardgrid}{%
  \foreach \gl in {1,2} {\draw[gridline,line width=0.4pt] (\gl,0) -- (\gl,3) (0,\gl) -- (3,\gl);}%
  \draw[ink,line width=1.2pt] (0,0) rectangle (3,3);%
}
% Для пустой доски: сетка той же толщины и цвета, что и граница.
\newcommand{\boardgridbold}{%
  \foreach \gl in {1,2} {\draw[ink,line width=1.2pt] (\gl,0) -- (\gl,3) (0,\gl) -- (3,\gl);}%
  \draw[ink,line width=1.2pt] (0,0) rectangle (3,3);%
}

% =====================================================================================================================
% Подписи и заголовок для доски 5,4 см. #1 --- стиль подписи, #2 --- заголовка.
% Смещение 0,70 см пересчитано в единицы доски делением на размер клетки.
% =====================================================================================================================
\newcommand{\boardlabels}[1]{%
  \foreach \cl/\lab in {0/a,1/b,2/c} {%
    \node[#1] at (\cl+0.5,{-0.70/5.4}) {\lab};
    \node[#1] at (\cl+0.5,{3.0+0.70/5.4}) {\lab};}%
  \foreach \rw/\lab in {0/3,1/2,2/1} {%
    \node[#1] at ({-0.70/5.4},{2.5-\rw}) {\lab};
    \node[#1] at ({3.0+0.70/5.4},{2.5-\rw}) {\lab};}%
}
\newcommand{\boardtitle}[2]{%
  \node[#2] at (1.5,{3.0+3.15/5.4}) {#1};%
}

% =====================================================================================================================
% Пешки по номерам квадратов. Радиус задан в единицах доски.
% =====================================================================================================================
\newcommand{\drawpawns}[2]{%
  \foreach \sqr in {#2} {\sqcx{\sqr}\sqcy{\sqr}\fill[pawn] (\qx,\qy) circle[radius=0.34];}%
  \foreach \sqr in {#1} {\sqcx{\sqr}\sqcy{\sqr}%
    \filldraw[fill=white,draw=pawn,line width=1.1pt] (\qx,\qy) circle[radius=0.34];}%
}

% =====================================================================================================================
% Небольшие иллюстрации к тексту правил и примеров (доска 1,9 см).
% Здесь пешки рисуются кружками-узлами, а не кругами: на маленьком размере так они лучше читаются.
% =====================================================================================================================
\tikzset{
  rulesheet/.style={
    x=1.9cm, y=1.9cm,
    line width=0.5pt, baseline=(current bounding box.center)},
  plainpawn/.style={circle, draw=pawn, line width=0.8pt, minimum size=7mm, inner sep=0pt},
  blpawn/.style={plainpawn, fill=pawn, draw=pawn},
  whpawn/.style={plainpawn, fill=white, draw=pawn},
  moveline/.style={-{Stealth[length=2.4mm,width=1.7mm]}, line width=0.9pt, draw=moveink},
}
\newcommand{\rulesheetgrid}{%
  \foreach \gl in {1,2} {\draw[ink!35] (\gl,0) -- (\gl,3) (0,\gl) -- (3,\gl);}%
  \draw[ink, line width=1pt] (0,0) rectangle (3,3);%
}
% \board{<белые>}{<чёрные>}{подписи: 1 да, 0 нет}
\newcommand{\board}[3]{%
  \begin{tikzpicture}[rulesheet]
    \rulesheetgrid
    \foreach \s in {#1} {\sqcx{\s}\sqcy{\s}\node[whpawn] at (\qx,\qy) {};}
    \foreach \s in {#2} {\sqcx{\s}\sqcy{\s}\node[blpawn] at (\qx,\qy) {};}
    \ifnum#3=1\relax
      \foreach \c/\l in {0/a,1/b,2/c} {%
        \node[font=\scriptsize\sffamily, text=ink!70] at (\c+0.5,{-0.34}) {\l};
        \node[font=\scriptsize\sffamily, text=ink!70] at (\c+0.5,{3.34}) {\l};}
      \foreach \r in {1,2,3} {%
        \node[font=\scriptsize\sffamily, text=ink!70] at ({-0.32},{\r-0.5}) {\r};
        \node[font=\scriptsize\sffamily, text=ink!70] at ({3.32},{\r-0.5}) {\r};}
    \fi
  \end{tikzpicture}%
}
% \boardmoves{<белые>}{<чёрные>}{<ходы: from/to, ...>}
\newcommand{\boardmoves}[3]{%
  \begin{tikzpicture}[rulesheet]
    \rulesheetgrid
    \foreach \s in {#1} {\sqcx{\s}\sqcy{\s}\node[whpawn] at (\qx,\qy) {};}
    \foreach \s in {#2} {\sqcx{\s}\sqcy{\s}\node[blpawn] at (\qx,\qy) {};}
    \foreach \f/\t in {#3} {%
      \sqcx{\f}\sqcy{\f}\pgfmathsetmacro{\fx}{\qx}\pgfmathsetmacro{\fy}{\qy}%
      \sqcx{\t}\sqcy{\t}%
      \draw[moveline] (\fx,\fy) -- (\qx,\qy);}
  \end{tikzpicture}%
}
% Крестик взятия: \take в тексте.
\newcommand{\take}{\ensuremath{\times}}

% =====================================================================================================================
% Карточка целиком (доска 5,4 см). Здесь пешки — круги, стрелки обрезаются
% по краям, чтобы пешки оставались видимыми.
% =====================================================================================================================
\tikzset{
  cardsheet/.style={x=5.4cm, y=5.4cm, line width=0.7pt, baseline=(current bounding box.center)},
  movetip/.style={-{Stealth[length=4mm,width=3.2mm]}, line width=5pt},
  squarelab/.style={font=\sffamily\bfseries\fontsize{36}{44}\selectfont, text=ink},
  movelab/.style={font=\sffamily\bfseries\fontsize{28}{34}\selectfont, text=ink},
  fieldlab/.style={font=\sffamily\bfseries\fontsize{36}{44}\selectfont,text=ink},
  fieldtitle/.style={font=\sffamily\fontsize{28}{34}\selectfont\bfseries,text=ink},
}
% Насколько конец стрелки не доходит до центра квадрата назначения:
% 0.13 для тихого хода, 0.54 для взятия.
\newcommand{\quietstop}{0.13}
\newcommand{\capturestop}{0.54}
% \card{<белые>}{<чёрные>}{<ходы>}{<подпись>}
% ходы перечисляются как <откуда>/<куда>/<цвет>, например 3/6/arrow1 по одному на каждый допустимый ход из этой позиции.
% Стрелка идёт от пешки к пешке, но обрезается по краям:
%   начало выходит из-под своей пешки,
%   конец уходит под пешку назначения (тихий ход) или останавливается перед чужой (взятие).
% Так пешки всегда остаются видимыми.
\newcommand{\card}[4]{%
  \begin{tikzpicture}[cardsheet]
    \boardgrid
    \foreach \mvfrom/\mvto/\mvcolor in {#3} {%
      \global\let\stopat\quietstop
      \foreach \sqr in {#1} {\ifnum\sqr=\mvto\relax\global\let\stopat\capturestop\fi}%
      \foreach \sqr in {#2} {\ifnum\sqr=\mvto\relax\global\let\stopat\capturestop\fi}%
      \sqcx{\mvfrom}\sqcy{\mvfrom}\pgfmathsetmacro{\fx}{\qx}\pgfmathsetmacro{\fy}{\qy}%
      \sqcx{\mvto}\sqcy{\mvto}%
      \pgfmathsetmacro{\ux}{\qx-\fx}\pgfmathsetmacro{\uy}{\qy-\fy}%
      \pgfmathsetmacro{\ulen}{sqrt(\ux*\ux+\uy*\uy)}%
      \pgfmathsetmacro{\sx}{round((\fx+0.40*\ux/\ulen)*100)/100}%
      \pgfmathsetmacro{\sy}{round((\fy+0.40*\uy/\ulen)*100)/100}%
      \pgfmathsetmacro{\ex}{round((\qx-\stopat*\ux/\ulen)*100)/100}%
      \pgfmathsetmacro{\ey}{round((\qy-\stopat*\uy/\ulen)*100)/100}%
      \draw[movetip,\mvcolor] (\sx,\sy) -- (\ex,\ey);%
    }%
    \drawpawns{#1}{#2}%
    \boardlabels{squarelab}%
    \boardtitle{#4}{movelab}%
  \end{tikzpicture}%
}
% Карточка целиком: доска по центру листа.
\newcommand{\hexcard}[4]{%
  \vspace*{\fill}%
  \begin{center}\card{#1}{#2}{#3}{#4}\end{center}%
  \vspace*{\fill}%
}
% Колонтитул: «страница/всего» слева, «карточка/колода» справа.
% Для страницы начальной расстановки номер в колоде не печатается.
\newcommand{\cardfoot}[2]{%
  \noindent\hbox to \textwidth{\color{black!30}\sffamily\fontsize{9}{11}\selectfont
    #1\if\relax\detokenize{#2}\relax\else\hfil #2\fi}%
  \vspace*{1.4cm}%
}
 из rules.tex,
% examples.tex, cards.tex, field.tex и guide.tex. rating.tex не подключает:
% ему нужны из палитры только цвета, поэтому он берёт colors.tex напрямую.
% Содержит только определения — \documentclass здесь быть не должно.
%
% Квадраты нумеруются слева направо и снизу вверх, как на настоящей доске:
%   7 8 9   <- ряд 3
%   4 5 6   <- ряд 2
%   1 2 3   <- ряд 1
% То есть a1=1, b1=2, c1=3, a2=4, …, c3=9.
% Ряд 1 внизу: белые стоят в ряду 1, чёрные в ряду 3.
% Подписи совпадают с подписями на игровом поле.
%
% Что где живёт:
%   \sqcx, \sqcy      — геометрия квадрата по его номеру;
%   \boardgrid        — тонкая сетка (карточки, игровое поле);
%   \boardgridbold    — сетка одной толщины с границей (пустое поле);
%   \boardlabels      — подписи a/b/c и 3/2/1 по краям доски 5,4 см;
%   \boardtitle       — заголовок над доской 5,4 см;
%   \drawpawns        — пешки, заданные номерами квадратов;
%   \board, \boardmoves — небольшие иллюстрации к тексту (доска 1,9 см);
%   \card, \hexcard, \cardfoot — карточка целиком (доска 5,4 см).
%
% Масштаб доски задаётся стилем TikZ, который выбирает вызывающий документ:
% rulesheet — 1,9 см (иллюстрации в тексте),
% cardsheet — 5,4 см (карточки).
% =====================================================================================================================

\RequirePackage{tikz}
\usetikzlibrary{arrows.meta}

% Палитра вынесена в colors.tex, см. там же список цветов и что ими красится.
% Общая палитра цветов для всего комплекта.
% Меняйте значения HEX и комментарии, имена менять не рекомендуется.
% Стрелки arrow1…arrow4 — условные названия, не переименовывать.

\RequirePackage{xcolor}

% Основная палитра для печати игровой доски и расстановки пешек
\definecolor{ink}{HTML}{26262A}       % текст, рамка доски, подписи клеток, оси в rating
\definecolor{pawn}{HTML}{000000}      % цвет пешек чёрных и контура пешек белых
\definecolor{gridline}{HTML}{B4B4BC}  % внутренние линии сетки
% Отдельный светлый цвет для стрелок на иллюстрациях: такая стрелка должна
% быть светлее чёрной пешки, иначе при взятии (например a2xb3) кончик
% стрелки пропадает на пешке. Поэтому у moveline свой цвет, а не ink.
\definecolor{moveink}{HTML}{B4B4BC}   % стрелки ходов в иллюстрациях и примерах

% Цвета стрелок: каждой стрелке на карточке соответствует конфета того же цвета.
\definecolor{arrow1}{HTML}{00B050}    % зелёный
\definecolor{arrow2}{HTML}{7030A0}    % фиолетовый
\definecolor{arrow3}{HTML}{FF0000}    % красный
\definecolor{arrow4}{HTML}{FFFF00}    % жёлтый



% =====================================================================================================================
% Геометрия квадрата по его номеру
% =====================================================================================================================
\newcommand{\sqcx}[1]{\pgfmathsetmacro{\qx}{mod(#1-1,3)+0.5}}       % центр по X
\newcommand{\sqcy}[1]{\pgfmathsetmacro{\qy}{floor((#1-1)/3)+0.5}}   % центр по Y

% =====================================================================================================================
% Сетка доски
% =====================================================================================================================
% Тонкая: внутренние линии светло-серые, граница толще. Карточки и поле.
\newcommand{\boardgrid}{%
  \foreach \gl in {1,2} {\draw[gridline,line width=0.4pt] (\gl,0) -- (\gl,3) (0,\gl) -- (3,\gl);}%
  \draw[ink,line width=1.2pt] (0,0) rectangle (3,3);%
}
% Для пустой доски: сетка той же толщины и цвета, что и граница.
\newcommand{\boardgridbold}{%
  \foreach \gl in {1,2} {\draw[ink,line width=1.2pt] (\gl,0) -- (\gl,3) (0,\gl) -- (3,\gl);}%
  \draw[ink,line width=1.2pt] (0,0) rectangle (3,3);%
}

% =====================================================================================================================
% Подписи и заголовок для доски 5,4 см. #1 --- стиль подписи, #2 --- заголовка.
% Смещение 0,70 см пересчитано в единицы доски делением на размер клетки.
% =====================================================================================================================
\newcommand{\boardlabels}[1]{%
  \foreach \cl/\lab in {0/a,1/b,2/c} {%
    \node[#1] at (\cl+0.5,{-0.70/5.4}) {\lab};
    \node[#1] at (\cl+0.5,{3.0+0.70/5.4}) {\lab};}%
  \foreach \rw/\lab in {0/3,1/2,2/1} {%
    \node[#1] at ({-0.70/5.4},{2.5-\rw}) {\lab};
    \node[#1] at ({3.0+0.70/5.4},{2.5-\rw}) {\lab};}%
}
\newcommand{\boardtitle}[2]{%
  \node[#2] at (1.5,{3.0+3.15/5.4}) {#1};%
}

% =====================================================================================================================
% Пешки по номерам квадратов. Радиус задан в единицах доски.
% =====================================================================================================================
\newcommand{\drawpawns}[2]{%
  \foreach \sqr in {#2} {\sqcx{\sqr}\sqcy{\sqr}\fill[pawn] (\qx,\qy) circle[radius=0.34];}%
  \foreach \sqr in {#1} {\sqcx{\sqr}\sqcy{\sqr}%
    \filldraw[fill=white,draw=pawn,line width=1.1pt] (\qx,\qy) circle[radius=0.34];}%
}

% =====================================================================================================================
% Небольшие иллюстрации к тексту правил и примеров (доска 1,9 см).
% Здесь пешки рисуются кружками-узлами, а не кругами: на маленьком размере так они лучше читаются.
% =====================================================================================================================
\tikzset{
  rulesheet/.style={
    x=1.9cm, y=1.9cm,
    line width=0.5pt, baseline=(current bounding box.center)},
  plainpawn/.style={circle, draw=pawn, line width=0.8pt, minimum size=7mm, inner sep=0pt},
  blpawn/.style={plainpawn, fill=pawn, draw=pawn},
  whpawn/.style={plainpawn, fill=white, draw=pawn},
  moveline/.style={-{Stealth[length=2.4mm,width=1.7mm]}, line width=0.9pt, draw=moveink},
}
\newcommand{\rulesheetgrid}{%
  \foreach \gl in {1,2} {\draw[ink!35] (\gl,0) -- (\gl,3) (0,\gl) -- (3,\gl);}%
  \draw[ink, line width=1pt] (0,0) rectangle (3,3);%
}
% \board{<белые>}{<чёрные>}{подписи: 1 да, 0 нет}
\newcommand{\board}[3]{%
  \begin{tikzpicture}[rulesheet]
    \rulesheetgrid
    \foreach \s in {#1} {\sqcx{\s}\sqcy{\s}\node[whpawn] at (\qx,\qy) {};}
    \foreach \s in {#2} {\sqcx{\s}\sqcy{\s}\node[blpawn] at (\qx,\qy) {};}
    \ifnum#3=1\relax
      \foreach \c/\l in {0/a,1/b,2/c} {%
        \node[font=\scriptsize\sffamily, text=ink!70] at (\c+0.5,{-0.34}) {\l};
        \node[font=\scriptsize\sffamily, text=ink!70] at (\c+0.5,{3.34}) {\l};}
      \foreach \r in {1,2,3} {%
        \node[font=\scriptsize\sffamily, text=ink!70] at ({-0.32},{\r-0.5}) {\r};
        \node[font=\scriptsize\sffamily, text=ink!70] at ({3.32},{\r-0.5}) {\r};}
    \fi
  \end{tikzpicture}%
}
% \boardmoves{<белые>}{<чёрные>}{<ходы: from/to, ...>}
\newcommand{\boardmoves}[3]{%
  \begin{tikzpicture}[rulesheet]
    \rulesheetgrid
    \foreach \s in {#1} {\sqcx{\s}\sqcy{\s}\node[whpawn] at (\qx,\qy) {};}
    \foreach \s in {#2} {\sqcx{\s}\sqcy{\s}\node[blpawn] at (\qx,\qy) {};}
    \foreach \f/\t in {#3} {%
      \sqcx{\f}\sqcy{\f}\pgfmathsetmacro{\fx}{\qx}\pgfmathsetmacro{\fy}{\qy}%
      \sqcx{\t}\sqcy{\t}%
      \draw[moveline] (\fx,\fy) -- (\qx,\qy);}
  \end{tikzpicture}%
}
% Крестик взятия: \take в тексте.
\newcommand{\take}{\ensuremath{\times}}

% =====================================================================================================================
% Карточка целиком (доска 5,4 см). Здесь пешки — круги, стрелки обрезаются
% по краям, чтобы пешки оставались видимыми.
% =====================================================================================================================
\tikzset{
  cardsheet/.style={x=5.4cm, y=5.4cm, line width=0.7pt, baseline=(current bounding box.center)},
  movetip/.style={-{Stealth[length=4mm,width=3.2mm]}, line width=5pt},
  squarelab/.style={font=\sffamily\bfseries\fontsize{36}{44}\selectfont, text=ink},
  movelab/.style={font=\sffamily\bfseries\fontsize{28}{34}\selectfont, text=ink},
  fieldlab/.style={font=\sffamily\bfseries\fontsize{36}{44}\selectfont,text=ink},
  fieldtitle/.style={font=\sffamily\fontsize{28}{34}\selectfont\bfseries,text=ink},
}
% Насколько конец стрелки не доходит до центра квадрата назначения:
% 0.13 для тихого хода, 0.54 для взятия.
\newcommand{\quietstop}{0.13}
\newcommand{\capturestop}{0.54}
% \card{<белые>}{<чёрные>}{<ходы>}{<подпись>}
% ходы перечисляются как <откуда>/<куда>/<цвет>, например 3/6/arrow1 по одному на каждый допустимый ход из этой позиции.
% Стрелка идёт от пешки к пешке, но обрезается по краям:
%   начало выходит из-под своей пешки,
%   конец уходит под пешку назначения (тихий ход) или останавливается перед чужой (взятие).
% Так пешки всегда остаются видимыми.
\newcommand{\card}[4]{%
  \begin{tikzpicture}[cardsheet]
    \boardgrid
    \foreach \mvfrom/\mvto/\mvcolor in {#3} {%
      \global\let\stopat\quietstop
      \foreach \sqr in {#1} {\ifnum\sqr=\mvto\relax\global\let\stopat\capturestop\fi}%
      \foreach \sqr in {#2} {\ifnum\sqr=\mvto\relax\global\let\stopat\capturestop\fi}%
      \sqcx{\mvfrom}\sqcy{\mvfrom}\pgfmathsetmacro{\fx}{\qx}\pgfmathsetmacro{\fy}{\qy}%
      \sqcx{\mvto}\sqcy{\mvto}%
      \pgfmathsetmacro{\ux}{\qx-\fx}\pgfmathsetmacro{\uy}{\qy-\fy}%
      \pgfmathsetmacro{\ulen}{sqrt(\ux*\ux+\uy*\uy)}%
      \pgfmathsetmacro{\sx}{round((\fx+0.40*\ux/\ulen)*100)/100}%
      \pgfmathsetmacro{\sy}{round((\fy+0.40*\uy/\ulen)*100)/100}%
      \pgfmathsetmacro{\ex}{round((\qx-\stopat*\ux/\ulen)*100)/100}%
      \pgfmathsetmacro{\ey}{round((\qy-\stopat*\uy/\ulen)*100)/100}%
      \draw[movetip,\mvcolor] (\sx,\sy) -- (\ex,\ey);%
    }%
    \drawpawns{#1}{#2}%
    \boardlabels{squarelab}%
    \boardtitle{#4}{movelab}%
  \end{tikzpicture}%
}
% Карточка целиком: доска по центру листа.
\newcommand{\hexcard}[4]{%
  \vspace*{\fill}%
  \begin{center}\card{#1}{#2}{#3}{#4}\end{center}%
  \vspace*{\fill}%
}
% Колонтитул: «страница/всего» слева, «карточка/колода» справа.
% Для страницы начальной расстановки номер в колоде не печатается.
\newcommand{\cardfoot}[2]{%
  \noindent\hbox to \textwidth{\color{black!30}\sffamily\fontsize{9}{11}\selectfont
    #1\if\relax\detokenize{#2}\relax\else\hfil #2\fi}%
  \vspace*{1.4cm}%
}
